% Real feature-importance bar chart (XGBoost, evaluation dataset, Section VI)
\begin{figure}[!t]
\centering
\begin{tikzpicture}[scale=1]
\foreach \name/\val [count=\i from 0] in {
    roll\_max\_14/0.0964,
    ewm\_7/0.0711,
    cos\_month/0.0682,
    lag\_21/0.0605,
    month/0.0600,
    lag\_1/0.0543,
    roll\_max\_7/0.0514,
    lag\_28/0.0505
} {
    \pgfmathsetmacro{\ypos}{-0.5*\i}
    \pgfmathsetmacro{\barlen}{\val*40}
    \node[anchor=east, font=\scriptsize] at (-0.15,\ypos) {\texttt{\name}};
    \draw[fill=blue!35, draw=blue!60] (0,\ypos-0.15) rectangle (\barlen/10,\ypos+0.15);
    \node[anchor=west, font=\tiny] at (\barlen/10+0.05,\ypos) {\val};
}
\draw[-{Latex[length=2mm]}] (0,-4.1) -- (4.2,-4.1) node[right, font=\tiny]{gain-based importance};
\end{tikzpicture}
\caption{Top eight gain-based feature importances from the XGBoost model fit on the evaluation dataset of Section~\ref{sec:dataset} (real values reported by the trained model, not illustrative data). Calendar features (\texttt{cos\_month}, \texttt{month}) and short/medium-range history (\texttt{lag\_1}, \texttt{lag\_21}, \texttt{lag\_28}, rolling maxima) dominate, consistent with the absence of strong exogenous signal in this particular dataset.}
\label{fig:feature-importance}
\end{figure}
